% exp-fips203-mlkem-kem-encaps-k3 — FIPS 203 Alg.20 ML-KEM-768 KEM Encaps
% 编译：bash scripts/xelatex-clean.sh examples/incubating/ml-kem/ml-kem-768/exp-fips203-mlkem-kem-encaps-k3/exp-fips203-mlkem-kem-encaps-k3-实现方案-customspec.tex
\documentclass[UTF8,a4paper,11pt]{ctexart}
\usepackage{amsmath,amssymb}
\usepackage{booktabs}
\usepackage{geometry}
\usepackage{hyperref}
\usepackage{longtable}
\usepackage{array}
\usepackage{seqsplit}
\geometry{margin=2.2cm}
\newcolumntype{L}[1]{>{\raggedright\arraybackslash}p{#1}}
\newcommand{\apiname}[1]{\texttt{\seqsplit{#1}}}
\newcommand{\dirname}[1]{\texttt{\seqsplit{#1}}}
\hypersetup{colorlinks=true,linkcolor=blue,urlcolor=blue}

\title{exp-fips203-mlkem-kem-encaps-k3\\FIPS 203 Alg.20 ML-KEM-768 KEM Encaps 实现方案}
\author{ascendc 工程（ML-KEM-768 W4b / E20）}
\date{2026-07-26}

\begin{document}
\maketitle

\begin{abstract}
本文档描述 \dirname{examples/incubating/ml-kem/ml-kem-768/exp-fips203-mlkem-kem-encaps-k3/}：
在 \textbf{Atlas A2 / Ascend910B4} 上以 AscendC 实现 FIPS 203 \textbf{Algorithm 20}，即 ML-KEM-768（$k=3$）\textbf{KEM Encaps}。

\textbf{参数锁定}：几何、I/O 与 launch 全部锁定到
\dirname{docs/specs/fips203-mlkem768-parameter-card.md} §3.3 D20：
\texttt{ek=1184B}，\texttt{m=32B}，输出 \texttt{c=1088B} 与 \texttt{K=32B}，SIM/生产 \textbf{2 launch}。
\textbf{行为基线}：活跃探针
\dirname{ascendc-tests/ml-kem/ml-kem-768/pass-fix-f203-alg20-kem-encaps-device-k3/}；
PKE Encrypt 主体可由本参数组 E14
\dirname{examples/incubating/ml-kem/ml-kem-768/exp-fips203-mlkem-pke-encrypt-k3/}
与 Alg.17 KEM 头组合。
\textbf{计划 AI Core 数}：Launch 1 为 AIV\_ONLY \texttt{blockDim=2}（KEM 头 + Encrypt prep），Launch 2 为 MIX \texttt{blockDim=1}（1 AIC + 2 AIV）。
\textbf{禁止}：创建 \dirname{examples/stable/ml-kem/ml-kem-768/}；从任何 \dirname{frozen/} 目录复制、改写或移植源码；沿用 k4 的 \texttt{c=1568}、\texttt{du=11/dv=5}、\texttt{re[9]} 或 pad-to-4/8 几何。
\end{abstract}

\tableofcontents
\newpage

\section{工程边界}

\subsection{允许的代码来源与依赖}

\begin{longtable}{@{}L{4.8cm}L{8.4cm}@{}}
\toprule
来源 & 用途 \\
\midrule
\dirname{ascendc-tests/ml-kem/ml-kem-768/pass-fix-f203-alg20-kem-encaps-device-k3/} & 活跃 D20 k3 行为基线；可复制后自包含适配到本 exp，最终不得保留对探针目录的编译期 include/import \\
\dirname{examples/incubating/ml-kem/ml-kem-768/exp-fips203-mlkem-pke-encrypt-k3/} & 活跃 E14 incubating PKE Encrypt；提供已自包含的 k3 Alg.14 主体 \\
\dirname{examples/incubating/ml-kem/ml-kem-1024/exp-fips203-mlkem-kem-encaps-k4/} & 活跃 k4 incubating KEM Encaps 结构参考；仅用于 retarget 到 k3 锁定字节长与几何 \\
\dirname{library/shared/} & SHAKE/Keccak、统一整数压缩、公共设备侧头文件等共享库 \\
CANN / tikicpulib / CAModel & 编译、CPU 孪生与 SIM 运行环境 \\
\bottomrule
\end{longtable}

\subsection{禁止项}

\begin{itemize}
  \item 禁止从 \dirname{ascendc-tests/frozen/} 或 \dirname{examples/frozen/} 拿源码、customspec 或路线。
  \item 禁止把 $h=H(\mathrm{ek})$、$r$、$y$、$e_1$、$e_2$、$\hat A$ 等中间态作为默认 I/O 落盘；debug dump 只能由显式开关启用。
  \item 禁止 Host 预填 $r$ 冒充设备 KEM 头；$r$ 必须由设备侧 $G(m\Vert H(ek))$ 产生并供 Encrypt prep 消费。
  \item 禁止为 $H$/$G$ 单独增加第三次生产 launch。
\end{itemize}

\section{数学语义（Alg.20 $\to$ Alg.17 $\to$ Alg.14）}

\subsection{输入输出}

\begin{longtable}{@{}L{3.5cm}L{9.8cm}@{}}
\toprule
张量 / 文件 & 契约 \\
\midrule
\texttt{input/ek\_kem.bin} & 1184B，等于 PKE public key；前 1152B 为 \(\hat t\) 的 \(\mathrm{ByteEncode}_{12}\)，末 32B 为 \(\rho\) \\
\texttt{input/m.bin} & 32B，Alg.20 internal message；设备侧 $G(m\Vert H(ek))$ 的前半为共享秘密 $K$，后半为 Encrypt 随机性 $r$ \\
\texttt{input/lut\_even\_stacked.bin}、\texttt{lut\_odd\_stacked.bin} & NTT/INTT Stage2 LUT；由本目录 golden 脚本生成；不是用户可见交付输入 \\
\texttt{output/c.bin} & 1088B，\(c=c_1\Vert c_2\)，其中 \(c_1=960\)B（\(d_u=10\)，3 poly），\(c_2=128\)B（\(d_v=4\)，1 poly） \\
\texttt{output/K.bin} & 32B，\(G(m\Vert H(ek))\) 前半；由设备侧 KEM 头写出 \\
\bottomrule
\end{longtable}

\subsection{阶段公式}

\begin{enumerate}
  \item KEM 头：$h=H(\mathrm{ek})=\mathrm{SHA3\textrm{-}256}(\mathrm{ek})$，$(K\Vert r)=G(m\Vert h)=\mathrm{SHA3\textrm{-}512}(m\Vert h)$。
  \item Alg.14 行 3--7：取 \(\rho=\texttt{ek[1152:1184]}\)，对 \((j,p)\in\{0,1,2\}^2\) 生成 \(\hat A[j,p]\)，GM 形状锁定为 \texttt{[9,256] int32}。
  \item Alg.14 行 8--15：用设备产生的 \(r\) 采样 \(y\Vert e_1\Vert e_2\) 共 7 个 poly；GM 形状锁定为 \texttt{[7,256]}。
  \item Alg.14 行 16--21：计算 \(\hat r=\mathrm{NTT}(y)\)，得 \(\hat u\) 与 \(\widehat{tr}\)，再对 \(u_0,u_1,u_2,v\) 执行真 \textbf{INTT batch=4}。
  \item Alg.14 行 22：按 \(d_u=10,d_v=4\) pack 得 \(c=c_1\Vert c_2\)；KEM 头已写出的 \(K\) 与 \(c\) 共同作为输出。
\end{enumerate}

\section{AscendC 实现规划}

\subsection{Launch 与 AI Core 规模}

\begin{longtable}{@{}L{2.5cm}L{3.1cm}L{7.5cm}@{}}
\toprule
Launch & 核类型 / blockDim & 数据划分与理由 \\
\midrule
1 prep+head & AIV\_ONLY，\texttt{blockDim=2} & block0 先执行 KEM 头，产生 \(K\) 与 \(r\)；随后继承 E14/D14 prep：$\hat A$ 共 9 poly 按 5+4 分给两个 AIV，\(y\Vert e_1\Vert e_2\) 共 7 poly 生成到真实 GM。选择 2 AIV 是参数卡 §3.3 D20 对 D14 prep 的锁定继承。 \\
2 compute+pack & MIX，\texttt{blockDim=1} & 1 AIC + 2 AIV；复用 E14/D14 compute，NTT(\(y\)) 只处理真实 3 poly，INTT batch4 连续 2+2，tail pack 输出 `c=1088B`。 \\
\bottomrule
\end{longtable}

\subsection{GM 几何}

\begin{longtable}{@{}L{3.8cm}L{9.3cm}@{}}
\toprule
对象 & 锁定几何 \\
\midrule
\(\hat A\) & \texttt{[9,256] int32}；9 个真实 poly，禁止 pad 到 16 \\
\texttt{re} & \texttt{[7,256] int32}；\texttt{y[0..2] || e1[0..2] || e2[0]}，随机性来自设备 \(r\) \\
\(\hat y\) NTT & polyvec3；复用 B5/NTT 真实 poly 几何，不引入假 poly \\
INTT batch4 & S0 \texttt{[8,256] int8}，Stage2 \texttt{mat\_c [32,128] int32}；Cube 的 M 对齐垫只是硬件对齐，不表示假 poly \\
Tail pack & \(u\) 三 poly 用 \(d_u=10\) pack 为 960B；\(v\) 一 poly 用 \(d_v=4\) pack 为 128B \\
\bottomrule
\end{longtable}

\subsection{同步与流水}

KEM 头写出 \(K\) 与 \(r\) 后，prep 的 PRF/CBD 必须在同一 launch 内读到完整 \(r\)；需要用本地已验的 AIV 内同步/顺序约束，禁止只凭 CPU 顺序假设 SIM 可见。
两次 launch 之间由 Host 顺序保证 GM 可见性。Compute 内部沿用 E14/D14 已验同步：AIV 写平面 Stage1/Stage3 GM 后以 CrossCore flag 与 AIC 交接；AIC 写完 \texttt{mat\_c} 后 AIV 等待再做 Stage3 merge、INTT 与 tail pack。
CPU 仅作为逻辑孪生；SIM 是同步与搬运验收门槛。

\section{AscendC API 表}

本实现不引入 E20 独有的新 AscendC 矢量/矩阵 API；KEM 头使用 shared Keccak `Sha3OneShot`。下表 API 均已有索引记录，使用约束按
\dirname{library/documents/CANN-AscendC算子开发接口参考-查阅索引.md} 执行。

\begin{longtable}{@{}L{3.2cm}L{2.7cm}L{2.4cm}L{2.3cm}L{3.0cm}@{}}
\toprule
API 名 & 分类 / 文档 & Atlas A2 & 本用例 dtype & 备注 \\
\midrule
\apiname{GetBlockIdx} / \apiname{GetSubBlockIdx} & 系统变量 / 资源 & √ & 标量索引 & 区分 AIV 分片与 MIX 子核；KEM 头仅固定 block 执行 \\
\apiname{GlobalTensor} / \apiname{LocalTensor} & 基础结构 & √ & \texttt{uint8/int8/int32} & GM/UB 张量契约必须与本 spec 几何一致 \\
\apiname{DataCopy} & Memory 数据搬运 & √ & \texttt{uint8/int8/int32} & 连续块搬运；长度与地址遵守 32B 对齐约束 \\
\apiname{Duplicate} & Memory 矢量计算 & √ & \texttt{uint8/int32} & 用于 UB 清零、pad 与 tail buffer 初始化 \\
\apiname{PipeBarrier} & Pipe 同步 & √ & — & MTE/Vector/Cube 阶段交界显式同步 \\
\apiname{CrossCoreSetFlag} / \apiname{CrossCoreWaitFlag} & MIX 同步 & √ & — & AIV/AIC 通过 GM 交接时使用；CPU 过不能删 \\
\apiname{Mmad} / \apiname{Fixpipe} & 矩阵计算 & √ & \texttt{int8*int8->int32} & NTT/INTT Stage2 MMAD；INTT batch4 的逻辑 M=8，硬件 M 对齐垫到 16 \\
\apiname{ShiftRight} & 矢量算术 & √ & \texttt{int32} & limb 拆分、Barrett 压缩与 ByteEncode 辅助右移 \\
\apiname{Muls} / \apiname{Add} / \apiname{Sub} & 矢量算术 & √ & \texttt{int32/int16} & limb、模加减、压缩/编码辅助 \\
\apiname{Cast} & 类型转换 & √（受限） & \texttt{int32/int16/half/int8} & 遵守索引中 A2 Cast 链约束；不得假设存在 \texttt{int32->int8} 单跳 \\
\apiname{Gather} & 矢量搬运/选择 & √ & \texttt{int32} & 仅允许在 Alg.11 / ByteEncode 等非 NTT S1--S3 段按既有路径使用；NTT 三段内禁用 \\
\apiname{GetValue} / \apiname{SetValue} & LocalTensor 标量访问 & √ & \texttt{uint8/int32} & 仅用于 pack 与 KEM 头尾标量缺口；不得把 Compares bit 包当逐 lane 布尔 \\
\bottomrule
\end{longtable}

\subsection{评估过但未采用}

\begin{longtable}{@{}L{4.0cm}L{9.0cm}@{}}
\toprule
API / 路线 & 不采用原因 \\
\midrule
\apiname{Compares}/\apiname{GatherMask} compact 链 & Alg.20 E20 不新增 reject compact；沿用 E14/D14 已验 pack 路径，不新增比较掩码风险 \\
Host 预填 `coins.bin` 作生产随机性 & Alg.20/17 要求 \(r=G(m\Vert H(ek))\)；Host `coins` 只能作为 golden 或调试对照，不作默认生产输入 \\
第三 launch KEM head & 参数卡 §3.3 D20 锁定 SIM 2 launch；KEM 头必须并入 prep launch \\
\bottomrule
\end{longtable}

\section{数学步骤覆盖率}

\begin{longtable}{@{}L{4.4cm}L{4.0cm}L{2.0cm}L{2.8cm}@{}}
\toprule
数学步骤 & 实现方式 / API & 覆盖 & 缺口说明 \\
\midrule
$H(ek)$ 与 $G(m\Vert H(ek))$ & shared Keccak SHA3-256/512 + UB/GM 拼接 & 100\% 设备 & `K` 与 `r` 均由设备 KEM 头产生 \\
取 \(\rho\)、采样 \(\hat A[9]\) & AIV\_ONLY + XOF/rej + \apiname{DataCopy} & 100\% 设备 & 5+4 分片；无零垫 \\
采样 \(y\Vert e_1\Vert e_2[7]\) & AIV\_ONLY + PRF/CBD & 100\% 设备 & 真实 7 poly；不借 k4 9 poly \\
NTT(\(y\)) 与内积 & MIX AIV/AIC + \apiname{Mmad} & 100\% 设备 & NTT 内无 \apiname{Gather} \\
INTT batch4 & MIX AIV/AIC + Stage3 merge & 100\% 设备 & 真 polyvec4，AIV 连续 2+2 \\
Compress/ByteEncode \(d=10/4\) & 统一整数压缩 + pack & 部分向量 & pack 尾部标量缺口沿用 E14/D14 已验实现 \\
\bottomrule
\end{longtable}

\section{验证计划}

\subsection{默认验收}

\begin{verbatim}
cd examples/incubating/ml-kem/ml-kem-768/exp-fips203-mlkem-kem-encaps-k3
bash run.sh -r cpu -v Ascend910B4
SIM_DIRECT=1 bash run.sh -r sim -v Ascend910B4
\end{verbatim}

期望：\texttt{c.bin} 为 1088B，\texttt{K.bin} 为 32B；两者均与 golden 对拍 PASS。
SIM 通过后记录 \texttt{Total tick} 到 \texttt{qa/active\_sim\_regress\_summary.md}，并检查用例根目录无 stray \texttt{core*.dump}、\texttt{profile\_*log*.toml}。

\subsection{可选对照}

后续可补 liboqs-768 encaps KAT；E20 本轮门禁为 CPU + \texttt{SIM\_DIRECT=1} sim。
liboqs 仅作为黑盒 I/O oracle；禁止把 liboqs 源码实现同构移植进 AscendC。

\section{产出物}

\begin{itemize}
  \item \dirname{exp-fips203-mlkem-kem-encaps-k3-实现方案-customspec.tex}
  \item \dirname{exp-fips203-mlkem-kem-encaps-k3-实现方案-customspec.pdf}
  \item 后续实现文件须与本 customspec 同目录自包含，且自研 Python/C/AscendC 写详细中文注释。
\end{itemize}

\end{document}
